The current expansion rate of the Universe, quantified by the Hubble constant ($H_0$), is a cornerstone of modern cosmology \citep[see reviews by][and references therein]{1992PASP..104..599J,2010ARA&A..48..673F}. Over the past decade, increasingly precise measurements of $H_0$ have revealed a striking and persistent discrepancy between its value inferred from observations of the early Universe, such as the cosmic microwave background (CMB), and its value measured directly in the local Universe using distance ladder methods. This disagreement, known as the ``Hubble Tension'', has persisted for a decade, exceeds the threshold for a statistical fluctuation, and has withstood extensive scrutiny of both observational data and analysis techniques \citep[see][for recent reviews]{2019NatAs...3..891V, 2021A&ARv..29....9S,2023ARNPS..73..153K,Verde2024}. 
As such, it poses a major challenge to the standard $\Lambda$CDM cosmological model and may point to new physics \citep[see][]{2021CQGra..38o3001D,CosmoVerse:2025txj}, barring the increasingly observationally disfavored possibility of multiple, independent, and unrecognized systematics. While many distance indicators have been used to measure the local value of {\Hcst}, few of these studies attempt to optimally combine these measures, which would require properly accounting for their correlations. In addition, correlations and/or redundant information offer a key advantage to ensure robustness.  Rather than a single distance ladder, or several parallel, partially correlated ladders, we show below that these measurements constitute a (stable) network: the Local Distance Network.  The study presented here is a comprehensive, community-wide effort in 2025 to construct this Local Distance Network via a broad, collaborative effort.


\subsection{The ISSI Bern Workshop setup and philosophy\label{sec:philosophy}}

Building a Local Distance Network requires expert knowledge across diverse astronomical disciplines. There is a wide range of distance indicators with varying levels of maturity, confidence, and uncertainty, necessitating careful consideration before employing them jointly. The various subsets of the astronomical community working on aspects related to the Hubble Tension have interacted at different junctures and have a general understanding of each other's methodologies. However, accurately and reliably combining results while considering all inter-dependencies requires a hands-on collaborative approach and a careful and thorough treatment, rooted in transparency, engagement, and scientific discourse. 

This was the underlying motivation and raison d'\^etre for the workshop\footnote{\url{https://workshops.issibern.ch/hubble-constant/}} ``What's under the H0od?'' held at the International Space Sciences Institute (ISSI) in Bern, Switzerland, in March 2025. The goal of this workshop was to arrive at a consensus set of `baseline' and `variant' datasets to include, to define statistically rigorous analysis procedures that account for dataset covariance, and to begin developing the open access tools required to measure {\Hcst} within a networked formalism. The Local Distance Network (Fig.~\ref{fig:Metro}) extends the distance ladder concept ``horizontally'' by linking multiple, overlapping calibration paths. It combines the statistical advantages of consistently averaging the contributions from multiple probes with the robustness to allow for the omission of any single probe. This is possible because there are multiple indicators that can serve the same methodological role (e.g., different anchors, different intermediate calibrators, different tracers of the Hubble flow, cf. Sect.\,\ref{sec:nomenclature}), as well as some parts of the distance network that require fewer connections or steps (megamaser distances, type II supernova modeling).  This ambitious program required assembling the leading experts in each of the relevant reserach areas.
The workshop conveners sought to leverage worldwide expertise across multiple tools as much as possible in selecting the $\sim$40 in-person attendees invited to participate in the ISSI workshop (see Section~\ref{expertise} for a list of attendees and their respective areas of expertise). Attention was given to inviting representatives from the most active groups in the field, including competing groups using similar or identical methods, in order to build consensus on how to consistently incorporate methods within the network and evaluate the level of agreement, especially in cases considered contentious in the literature.\footnote{Though the overall acceptance rate was high ($>$ 90\%), several invitees were unable to attend or declined participation in the workshop. Following the inclusive spirit of the workshop, opportunities for remote participation via Zoom and asynchronous participation by e-mail were offered in such cases. In the case where a group declined to participate in any format, their published data were still included to preserve their scientific contributions.} The group thus assembled (i.e., we, the H${}_{0}$DN collaboration) placed emphasis on the methodology of combining datasets (i.e., the ``how-to''), rather than on the results themselves, which were understood to be subject to quality and consistency checks (e.g., $\chi^2$ or residuals) and to be published irrespective of the resulting value of \Hcst.

\begin{figure*}
    \centering
    \includegraphics[]{Figures/LocalDistanceNetwork_Conceptual_02.pdf}
    \caption{Conceptual overview of the Local Distance Network, a many routes approach. Different methods for distance determination may connect the absolute scale determined by geometric means to \Hcst. A non-exhaustive list of baseline linkages discussed in the literature or the paper is labeled on the right. Links to geometric distances provided by Masers, DEB, and Parallax are indicated as available in our analysis. Background rectangles in orange, light blue, and gray indicate where Rung 1, Rung 2, and Rung 3 of a traditional distance ladder would fall. Unlabeled tickmarks represent Groups (Fornax \& Virgo for the TRGB to SBF, Coma for FP). Example references: 
    \citet[SH0ES]{2022ApJ...934L...7R}, \citet[CCHP]{2025ApJ...985..203F}, \citet[EDD]{2021AJ....162...80A}, \citet[Pop-II]{2024ApJ...973...83A}, \citet[MCP]{2020ApJ...890..118P}, \citet[Pop-I]{2022MNRAS.514.4620D}, \citet[CF4]{2020ApJ...902..145K}, \citet[DESI]{2025MNRAS.539.3627S}, \citet[adh0cc]{2025AA...702A..41V}. Appendix\,\ref{app:replication} replicates a subset of these routes.}
    \label{fig:Metro}
\end{figure*}



